\documentclass[11pt]{article}

\usepackage[margin=1in]{geometry}
\usepackage{amsmath,amssymb}
\usepackage{array,longtable}
\usepackage[hidelinks]{hyperref}

\setlength{\parindent}{0pt}
\setlength{\parskip}{0.6em}
\renewcommand{\arraystretch}{1.18}

\title{Two Clarifications: The BS20 Large-Economy LLN and Kline's
Wage-Difference Object}
\author{}
\date{}

\begin{document}
\maketitle

\section*{Answers in one paragraph}

In Borovi\v{c}kov\'a--Shimer (2020, henceforth BS20), \(I_x\) is the
baseline number of workers with characteristic \(x\), and \(\tau\) is a
theoretical replication index: the \(\tau\)-economy has \(\tau I_x\)
workers of characteristic \(x\) and \(\tau J_y\) firms of characteristic
\(y\).  The law of large numbers operates across the growing numbers of
workers and firms, not across an increasing number of matches for each
individual.  In Kline, the primitive wage-difference object
\(Y_{i2}(k)-Y_{i2}(j)\) is closest to Sorkin--Warwar's worker-specific firm
treatment effect; Volpe instead presents the two pillars through a wage-level
equation.  Kline's no-selection-on-treatment-effects restriction is
conceptually closer to exogenous mobility than to additive separability:
individual gains may remain heterogeneous.  Even when the average edge
effect can be written as a difference between average firm wage levels, that
does not imply, and is not equivalent to, additive separability of
individual potential wages.

\section{BS20: what goes to infinity in the LLN?}

\subsection{Notation}

\begin{longtable}{p{0.20\textwidth}p{0.53\textwidth}p{0.20\textwidth}}
\textbf{Symbol} & \textbf{Meaning} & \textbf{Source or status} \\
\hline
\(\mathcal X,\mathcal Y\) & Finite sets of possible worker and firm
characteristics. & BS20, p.\ 18; calligraphic set notation is a
clarification. \\
\(x,y\) & A worker characteristic and a firm characteristic. & BS20,
p.\ 18. \\
\(I_x,J_y\) & Baseline numbers of workers with characteristic \(x\) and
firms with characteristic \(y\). & BS20, p.\ 18. \\
\(I,J\) & Baseline total numbers of workers and firms:
\(I=\sum_{x\in\mathcal X}I_x\) and
\(J=\sum_{y\in\mathcal Y}J_y\).  These are counts, not random identities. &
BS20, p.\ 18. \\
\(\tau\) & Positive-integer replication index.  The \(\tau\)-economy has
\(\tau I_x\) workers of characteristic \(x\) and \(\tau J_y\) firms of
characteristic \(y\). & BS20, p.\ 18. \\
\(i,j\) & Individual worker and firm indices. & BS20, pp.\ 18--19. \\
\(x_i\) & Characteristic of worker \(i\). & BS20, pp.\ 18--21. \\
\(M_i,N_j\) & Numbers of distinct matches observed for worker \(i\) and firm
\(j\). & BS20, pp.\ 18--22. \\
\(m,m'\) & Two distinct match indices for worker \(i\). & BS20,
pp.\ 19--22. \\
\(w^w_{i,m}\) & Worker-side average log wage in worker \(i\)'s match \(m\). &
BS20, p.\ 19. \\
\(\lambda(x_i)\) & Expected-wage type of a worker with characteristic
\(x_i\). & BS20, pp.\ 5, 20--22. \\
\(\varepsilon^w_{i,m}\) & Worker-side match-wage error in the auxiliary
equation \(w^w_{i,m}=\lambda(x_i)+\varepsilon^w_{i,m}\). & BS20,
pp.\ 21--22. \\
\(S_i\) & Cross-match product score for \(\lambda(x_i)^2\), defined in
(B.5). & Derived shorthand for BS20 equation (18), p.\ 22. \\
\(\eta_i\) & Score error \(S_i-\lambda(x_i)^2\). & Derived notation. \\
\(A_\tau\) & Cross-worker average of \(\eta_i\) in the \(\tau\)-economy. &
Derived notation. \\
\(C\) & A finite uniform upper bound on
\(\operatorname{Var}(\eta_i)\). & Maintained regularity condition used in
the derived LLN proof. \\
\(\delta\) & Arbitrary positive tolerance in Chebyshev's inequality. &
Standard proof notation. \\
\(T_i^w,\overline T_x^w\) & Total employment time observed for worker \(i\),
and its characteristic-\(x\) expectation
\(\overline T_x^w=\mathbb E[T_i^w\mid x_i=x]\). & BS20, pp.\ 19--20;
conditional-mean notation is derived. \\
\(\widehat Q_{\lambda,\tau},Q_\lambda\) & Duration-weighted estimator and
population target for the second moment of worker wage types, defined in
(B.11)--(B.12). & Derived labels for BS20 equation (20), p.\ 23. \\
\(\mu(y)\) & Expected-wage type of a firm with characteristic \(y\). &
BS20, pp.\ 5, 20--23. \\
\(\widehat\rho_{\mathrm{BS}},\rho_{\mathrm{BS}}\) & BS20 sorting-correlation
estimator and population estimand. & BS20, pp.\ 5--6, 24. \\
\end{longtable}

\subsection{Real-world interpretation of the LLN variables}

The following table translates the statistical notation into the objects one
would recognize in matched worker--firm data.  The translations explain the
asymptotic experiment; they do not claim that every characteristic is
directly observed by the econometrician.

\begin{longtable}{p{0.16\textwidth}p{0.37\textwidth}p{0.37\textwidth}}
\textbf{Symbol} & \textbf{Actual-economy interpretation} &
\textbf{What it is not} \\
\hline
\(x\) & A complete worker characteristic in the statistical model: the
worker's wage-relevant, matching-relevant, and duration-relevant state.
Depending on the application, it may summarize observed attributes, latent
skill, or both. & It need not be a single observed label such as education,
and BS20's estimator does not require the econometrician to assign every
worker to an observed \(x\)-cell. \\
\(y\) & The analogous complete firm characteristic governing the firm's
matching, wage, employment-duration, and observation distributions. & It
need not be an observed industry or firm-size category. \\
\(I_x\) & Number of distinct worker identities with characteristic \(x\) in
the baseline theoretical economy.  For example, if the baseline contains
2,000 workers with state \(x\), then \(I_x=2{,}000\). & It is not worker
\(i\)'s number of jobs, wage observations, or years in the panel. \\
\(J_y\) & Number of distinct firm identities with characteristic \(y\) in
the baseline theoretical economy. & It is not the number of workers employed
by one firm. \\
\(I,J\) & Total numbers of distinct workers and firms in the baseline
economy.  In administrative data, these correspond to counts of worker IDs
and firm IDs in the analysis population. & They are not numbers of
worker--firm matches or person-year observations. \\
\(\tau\) & A theoretical market-size multiplier.  Moving from \(\tau=1\)
to \(\tau=5\) multiplies every worker-characteristic and firm-characteristic
population count by five while preserving their shares and
characteristic-specific data-generating processes. & It is not observed in
the data, not calendar time, not five copies of each realized wage, and not
five times as many observations on the same worker. \\
\(\tau I_x,\tau J_y\) & Numbers of new worker and firm identities of each
characteristic in the enlarged economy.  These identities receive new match,
wage, and duration realizations from the same characteristic-specific
distributions. & They are not literal duplicates with identical employment
histories. \\
\(i,j\) & A particular worker ID and firm ID. & They are identities, not
characteristic categories.  Many different workers \(i\) may share the same
\(x_i=x\). \\
\(M_i\) & Number of distinct worker--firm matches in worker \(i\)'s observed
history; empirically, this is approximately the number of distinct employers
after repeated spells with the same employer are combined as required by the
design. & It need not grow with \(\tau\).  A worker can have only two
qualifying matches in every economy in the sequence. \\
\(N_j\) & Number of distinct workers matched with firm \(j\) in its observed
history. & It need not grow with \(\tau\); it is not the total number of
firms. \\
\(w^w_{i,m}\) & Average log wage paid in worker \(i\)'s \(m\)-th qualifying
worker--firm match. & It is not worker \(i\)'s latent type and is not
necessarily an annual wage observation. \\
\(\lambda(x_i)\) & Expected log wage received by workers with characteristic
\(x_i\), using the worker-side, duration-weighted match distribution.  It is
the population worker wage type. & It is not consistently estimated for
each individual merely because that individual has two jobs. \\
\(\varepsilon^w_{i,m}\) & Difference between the realized match-average log
wage and the worker characteristic's expected-wage component in the
auxiliary equation. & It is not assumed to disappear for a given worker as
the economy grows. \\
\(S_i\) & A cross-product of worker \(i\)'s distinct-match wages designed so
that its conditional expectation equals \(\lambda(x_i)^2\).  It is one
noisy worker-level contribution to the aggregate second moment. & It is not
a consistent estimate of \(\lambda(x_i)^2\) when \(M_i\) remains fixed. \\
\(\eta_i\) & The worker-level measurement error in that moment score:
\(S_i-\lambda(x_i)^2\). & It is not a new wage shock or economic match
component; it is estimation noise induced by having few matches. \\
\(A_\tau\) & Average of the score errors \(\eta_i\) across the
\(\tau I\) workers.  This is the noise that the cross-sectional LLN removes.
& It is not the average error across increasingly many jobs of one worker. \\
\(T_i^w\) & Total time worker \(i\) is observed employed, used to make the
sample aggregation reproduce the duration weighting in the population
estimand. & It is not the replication index and need not be common across
workers. \\
\(\overline T_x^w\) & Average total observed employment time for workers
with characteristic \(x\) in the statistical population. & It is a
population conditional mean, not an additional observation for each worker.
\\
\(\widehat Q_{\lambda,\tau}\) & Duration-weighted sample second moment
formed by averaging the noisy \(S_i\)'s across all workers in the
\(\tau\)-economy. & It is not an individual worker-type estimator. \\
\(Q_\lambda\) & Corresponding population second moment of worker wage types.
Together with the mean, it yields the worker-type variance. & It is not a
finite-sample realization. \\
\(\mu(y)\) & Firm-side expected-wage type: the expected log wage paid by a
firm of characteristic \(y\). & It is not the wage in one worker--firm match.
\\
\(C,\delta\) & A finite variance bound and an arbitrary probability-tolerance
threshold used in the Chebyshev proof. & They are mathematical proof devices,
not features of a labor market. \\
\(\rho_{\mathrm{BS}},\widehat\rho_{\mathrm{BS}}\) & Population correlation
between matched worker and firm expected-wage types, and the statistic
computed from the data to estimate it. & The first is not observed; the
second is not the latent type of any worker or firm. \\
\end{longtable}

For a concrete illustration, suppose the baseline theoretical economy has
\(I=10{,}000\) workers and \(J=1{,}000\) firms, including
\(I_x=2{,}000\) workers of characteristic \(x\).  The economy with
\(\tau=5\) has \(50{,}000\) workers, \(5{,}000\) firms, and
\(10{,}000\) workers of characteristic \(x\).  The additional workers are
new identities with new employment histories drawn from the same
characteristic-\(x\) distribution.  A worker may still have only
\(M_i=2\) qualifying matches.  What becomes large is the number of distinct
worker-level scores \(S_i\), not the amount of information used to construct
each score.

In an actual administrative dataset, the researcher observes only one
economy and implements the formulas at \(\tau=1\).  Consistency interprets a
large observed economy as one member of the hypothetical sequence
\(\tau=1,2,\ldots\).  BS20's Auxiliary Assumption 4 supplies the
cross-unit independence needed for independently centered score noise to
average out; aggregate, regional, industry, or other common shocks can
invalidate that approximation if they leave substantial dependence across
units.

\textit{Source: BS20, pp.\ 18--24, especially the dynamic statistical model,
Auxiliary Assumptions 1--4, and equations (18)--(25).  The numerical example
and the empirical-language translations are derived.}

\subsection{The replicated economy}

BS20 begins with fixed baseline characteristic counts
\[
I=\sum_{x\in\mathcal X}I_x,
\qquad
J=\sum_{y\in\mathcal Y}J_y.
\tag{B.1}
\]
It then defines a sequence of economies indexed by a positive integer
\(\tau\).  In the \(\tau\)-economy the characteristic-specific counts are
\[
\tau I_x\quad\text{workers of characteristic }x,
\qquad
\tau J_y\quad\text{firms of characteristic }y.
\tag{B.2}
\]
Consequently, the total numbers of workers and firms are
\[
\text{total workers}=\tau I,
\qquad
\text{total firms}=\tau J.
\tag{B.3}
\]
The characteristic composition is unchanged because
\[
\frac{\tau I_x}{\tau I}
=
\frac{I_x}{I},
\qquad
\frac{\tau J_y}{\tau J}
=
\frac{J_y}{J}.
\tag{B.4}
\]
Thus \(\tau\to\infty\) means that the numbers of worker and firm identities
of every characteristic grow proportionally.  It is not calendar time, the
number of periods, or the number of observations on one fixed individual.
In particular, \(I_x\) is a baseline cross-sectional characteristic-cell
count; it is neither worker \(i\)'s number of observations nor worker
\(i\)'s number of matches.

\textit{Source: BS20, p.\ 18.  Equations (B.1)--(B.4) spell out the
paper's replication statement.}

\subsection{Why a fixed number of matches per worker is enough}

BS20 permits \(M_i\) to remain small as \(\tau\to\infty\).  Consider the
worker-side auxiliary equation
\[
w^w_{i,m}
=
\lambda(x_i)+\varepsilon^w_{i,m},
\qquad
\mathbb E[\varepsilon^w_{i,m}\mid x_i]=0.
\tag{B.5}
\]
For two distinct matches \(m\neq m'\), the conditional independence
assumption gives
\[
\begin{aligned}
&\mathbb E[
w^w_{i,m}w^w_{i,m'}\mid x_i]\\
&\quad=
\mathbb E[
(\lambda(x_i)+\varepsilon^w_{i,m})
(\lambda(x_i)+\varepsilon^w_{i,m'})
\mid x_i]\\
&\quad=
\lambda(x_i)^2
+\lambda(x_i)
\mathbb E[
\varepsilon^w_{i,m}+\varepsilon^w_{i,m'}\mid x_i]
+\mathbb E[
\varepsilon^w_{i,m}\varepsilon^w_{i,m'}\mid x_i]\\
&\quad=
\lambda(x_i)^2.
\end{aligned}
\tag{B.6}
\]
The two linear error terms vanish by conditional mean zero.  Conditional
independence makes the last expectation equal the product of the two
conditional error means, which is also zero.

Average the ordered distinct-match products:
\[
S_i
=
\frac{1}{M_i(M_i-1)}
\sum_{m=1}^{M_i}
\sum_{\substack{m'=1\\m'\neq m}}^{M_i}
w^w_{i,m}w^w_{i,m'}.
\tag{B.7}
\]
Equation (B.6) implies
\[
\mathbb E[S_i\mid x_i]=\lambda(x_i)^2.
\tag{B.8}
\]
This is an unbiasedness statement, not an individual consistency statement.
If \(M_i=2\) in every replicated economy, \(S_i\) remains noisy for that
worker.  BS20 does not need \(S_i\overset p\longrightarrow\lambda(x_i)^2\)
for each \(i\).

\textit{Source: BS20, pp.\ 21--23, Auxiliary Assumption 1 and equation
(18).  The expansion in (B.6) is derived.}

\subsection{The LLN, line by line}

Define
\[
\eta_i=S_i-\lambda(x_i)^2,
\qquad
A_\tau=\frac{1}{\tau I}\sum_{i=1}^{\tau I}\eta_i.
\tag{B.9}
\]
By (B.8), \(\mathbb E[\eta_i\mid x_i]=0\).  Suppose, as in the role played
by BS20's cross-unit independence and finite-moment conditions, that the
\(\eta_i\)'s are independent across workers and
\(\operatorname{Var}(\eta_i)\leq C<\infty\).  Then
\[
\begin{aligned}
\operatorname{Var}(A_\tau)
&=
\operatorname{Var}\left(
\frac{1}{\tau I}\sum_{i=1}^{\tau I}\eta_i
\right)\\
&=
\frac{1}{(\tau I)^2}
\sum_{i=1}^{\tau I}\operatorname{Var}(\eta_i)\\
&\leq
\frac{\tau I\,C}{(\tau I)^2}\\
&=
\frac{C}{\tau I}
\longrightarrow 0.
\end{aligned}
\tag{B.10}
\]
The second equality uses cross-worker independence.  The inequality uses
the uniform variance bound.  Chebyshev's inequality now gives, for every
\(\delta>0\),
\[
\Pr(|A_\tau|>\delta)
\leq
\frac{C}{\tau I\delta^2}
\longrightarrow0.
\tag{B.11}
\]
Therefore \(A_\tau\overset p\longrightarrow0\).  Because the replicated
economy contains exactly \(\tau I_x\) workers of characteristic \(x\),
\[
\begin{aligned}
\frac{1}{\tau I}\sum_{i=1}^{\tau I}S_i
&=
\frac{1}{\tau I}\sum_{i=1}^{\tau I}\lambda(x_i)^2+A_\tau\\
&=
\sum_{x\in\mathcal X}\frac{\tau I_x}{\tau I}\lambda(x)^2+A_\tau\\
&\overset p\longrightarrow
\sum_{x\in\mathcal X}\frac{I_x}{I}\lambda(x)^2.
\end{aligned}
\tag{B.12}
\]
The individual score noise does not disappear within a worker; it disappears
after averaging across the growing number \(\tau I\) of workers.

BS20's implemented second-moment estimator is duration weighted.  In the
present notation it has the form
\[
\widehat Q_{\lambda,\tau}
=
\frac{
\sum_{i=1}^{\tau I}T_i^wS_i
}{
\sum_{i=1}^{\tau I}T_i^w
},
\tag{B.13}
\]
with population target
\[
Q_\lambda
=
\frac{
\sum_{x\in\mathcal X}
I_x\overline T_x^w\lambda(x)^2
}{
\sum_{x\in\mathcal X}
I_x\overline T_x^w
}.
\tag{B.14}
\]
Under BS20's duration-error restrictions, cross-unit independence, and
finite-moment conditions, apply the same argument separately to the
normalized numerator and denominator of (B.13).  Both converge to their
counterparts in (B.14); the denominator has a positive limit, so the ratio
converges to \(Q_\lambda\).  The firm second moment follows by the same
argument across \(\tau J\) firms.  The covariance estimator uses
leave-focal-match-out worker and firm wage averages; BS20 first proves that
their product is unbiased for the matched type product and then shows that
the variance of its aggregate average is proportional to \(1/\tau\).

The final logic is therefore
\[
\begin{array}{c}
\text{two conditionally independent matches per unit}\\
\Downarrow\\
\text{unbiased unit-level scores for the required latent moments}\\
\Downarrow\\
\text{averaging across }\tau I\text{ workers and }\tau J\text{ firms}\\
\Downarrow\\
\text{aggregate moment consistency as }\tau\to\infty\\
\Downarrow\\
\widehat\rho_{\mathrm{BS}}\overset p\longrightarrow\rho_{\mathrm{BS}}.
\end{array}
\tag{B.15}
\]
The correlation step also requires positive worker- and firm-type variances.
In the empirical formulas BS20 sets \(\tau=1\); the sequence
\(\tau\to\infty\) is the theoretical device used to establish consistency.

\textit{Source: BS20, pp.\ 22--24, Propositions 3--5; Appendix B.1,
pp.\ 48--50.  Equations (B.9)--(B.14) isolate the paper's LLN argument in
elementary form.}

\section{Kline's wage difference versus Sorkin and Volpe}

\subsection{Notation}

\begin{longtable}{p{0.20\textwidth}p{0.53\textwidth}p{0.20\textwidth}}
\textbf{Symbol} & \textbf{Meaning} & \textbf{Source or status} \\
\hline
\(i,j,k,t,s,T\) & Worker, origin firm, destination firm, outcome period,
generic conditioning period, and total number of periods. & Standard
indices used by Kline, Sorkin--Warwar, and Volpe. \\
\(D_{it}\) & Firm employing worker \(i\) in period \(t\). & Kline,
pp.\ 29--33. \\
\(\mathbf D_i\) & Two-period mobility path
\((D_{i1},D_{i2})\). & Derived shorthand. \\
\(Y_{it}(j)\) & Worker \(i\)'s potential log wage in period \(t\) at firm
\(j\). & Kline, pp.\ 29--33. \\
\(Y_{it}\) & Observed log wage
\(Y_{it}=Y_{it}(D_{it})\). & Standard consistency relation; Kline,
pp.\ 29--30. \\
\(X_{it},\beta\) & Worker-time covariates and their coefficient. &
Sorkin--Warwar, p.\ 5; Volpe, slide 4/35. \\
\(\alpha_i,\psi_j\) & Additive worker and firm wage components. &
Sorkin--Warwar, pp.\ 5--7; Volpe, slide 4/35; Kline, p.\ 15. \\
\(\varpi_{ij}\) & Permanent statistical worker--firm match component.  This
relabels Sorkin--Warwar's \(\mu_{ij}\) to avoid collision with other uses of
\(\mu\). & Sorkin--Warwar, p.\ 6; symbol relabeled. \\
\(\widetilde\varepsilon_{it}(j)\) & Transitory potential-wage remainder at
firm \(j\). & Sorkin--Warwar, pp.\ 6--7; potential-firm index made
explicit. \\
\(\varepsilon^A_{it}\) & Composite residual in the observed additive AKM
equation. & Derived from Sorkin--Warwar equation (2), p.\ 6. \\
\(d,\mathbf x\) & Generic realized firm and covariate values used in
mobility-conditioning cells. & Derived notation. \\
\(\mathsf E_{\widetilde\varepsilon}\) & Conditional-mean exogeneity of the
transitory potential-wage remainder, defined in (K.3). & Derived label for
Sorkin--Warwar equation (3), p.\ 6. \\
\(\mathsf E_{\varpi}\) & Conditional-mean exogeneity of the permanent match
component, defined in (K.4). & Derived label for Sorkin--Warwar equation
(4), p.\ 6. \\
\(\mathsf A_0\) & Pure additive-separability restriction
\(\varpi_{ij}=0\) for every potential worker--firm pair. & Derived label for
Sorkin--Warwar equation (6), p.\ 7, and Volpe's second pillar. \\
\(\mathsf{EM}_{SW},\mathsf{AS}_{SW}\) & Sorkin--Warwar's exogenous-mobility
and additive-separability assumptions. & Derived labels. \\
\(\mathsf{EM}_{V},\mathsf{AS}_{V}\) & Volpe's exogenous-mobility and pure
additive-separability pillars. & Derived labels. \\
\(G_{i2}(j,k)\) & Worker \(i\)'s period-2 potential-wage gain from replacing
firm \(j\) with firm \(k\). & Derived shorthand for Kline, pp.\ 30--33. \\
\(\mathsf{NSG}_{K}\) & Kline's no-selection-on-treatment-effects
restriction, defined in (K.9). & Kline, p.\ 32; label derived. \\
\(\Delta_{jk}\) & Average causal wage effect for workers selected onto the
directed edge \(j\to k\). & Kline, pp.\ 30--31. \\
\(\psi_j^{\mathrm{ATE}}\) & Population mean potential-wage level at firm
\(j\), relative to reference firm \(0\). & Derived relabeling of Kline,
p.\ 33. \\
\(\xi_i\) & Nondegenerate, mean-zero permanent match heterogeneity in the
two-firm counterexample. & Derived notation. \\
\end{longtable}

\subsection{One wage surface that separates the restrictions}

Use the potential-wage decomposition
\[
Y_{it}(j)
=
\alpha_i+\psi_j+X_{it}'\beta
+\varpi_{ij}
+\widetilde\varepsilon_{it}(j).
\tag{K.1}
\]
The observed wage is \(Y_{it}=Y_{it}(D_{it})\).  If an additive AKM equation
omits \(\varpi_{ij}\), its observed residual is
\[
\varepsilon^A_{it}
=
\varpi_{iD_{it}}
+\widetilde\varepsilon_{it}(D_{it}).
\tag{K.2}
\]
Equation (K.1) is a common bookkeeping device: it separates a permanent
worker--firm interaction from a transitory remainder.  It does not assume
that either component is exogenous.

\textit{Source: Sorkin--Warwar, pp.\ 5--7, equations (1)--(2).
The potential-wage notation and the relabeling
\(\mu_{ij}\mapsto\varpi_{ij}\) are derived clarifications.}

\subsection{Sorkin and Volpe use different packaging}

In the notation of (K.1), define the two componentwise mean-exogeneity
conditions, for every admissible \(i,j,d,s,t,\mathbf x\),
\[
\mathsf E_{\widetilde\varepsilon}:
\quad
\mathbb E[
\widetilde\varepsilon_{it}(j)
\mid D_{is}=d,X_{is}=\mathbf x]
=0,
\tag{K.3}
\]
\[
\mathsf E_{\varpi}:
\quad
\mathbb E[
\varpi_{ij}
\mid D_{is}=d,X_{is}=\mathbf x]
=0,
\tag{K.4}
\]
and define pure additive separability as
\[
\mathsf A_0:
\quad
\varpi_{ij}=0
\quad\text{for every }(i,j).
\tag{K.5}
\]

Sorkin--Warwar package the assumptions as
\[
\mathsf{EM}_{SW}
=
\mathsf E_{\widetilde\varepsilon}
\cap
\mathsf E_{\varpi},
\qquad
\mathsf{AS}_{SW}
=
\mathsf E_{\widetilde\varepsilon}
\cap
\mathsf A_0.
\tag{K.6}
\]
Because \(\mathsf A_0\) implies \(\mathsf E_{\varpi}\),
\[
\mathsf{AS}_{SW}\subseteq\mathsf{EM}_{SW}.
\tag{K.7}
\]
This is why Sorkin can describe additive separability as the stronger
constant-effects case and exogenous mobility as the weaker average-effects
case.  Sorkin's Lecture 10 summarizes the distinction as ``exogenous
mobility: average effects'' versus ``additive separability: constant
effects.''

Volpe presents two separate pillars:
\[
\mathsf{EM}_{V}:
\quad
\mathbb E\left[
\varepsilon^A_{it}
\mid
\alpha_i,\{\psi_{D_{is}},X_{is}\}_{s=1}^{T}
\right]=0
\quad\text{for every }t,
\qquad
\mathsf{AS}_{V}:
\quad
\varpi_{ij}=0
\quad\text{for every }(i,j).
\tag{K.8}
\]
As pure restrictions, a selection restriction and an outcome-surface
restriction are not nested.  The apparent disagreement is therefore about
packaging: Sorkin's named additive-separability assumption includes the
transitory exogeneity condition, while Volpe uses ``additive separability''
for the outcome restriction alone.  An AKM causal interpretation uses both
of Volpe's pillars.

\textit{Sources: Sorkin--Warwar, pp.\ 6--7, Assumptions 1--2; Sorkin
Lecture 10, slide 46/79; Volpe, slide 4/35.  The set notation in
(K.6)--(K.8) is derived.}

\subsection{Which framework uses the wage object closest to Kline's?}

Kline's period-2 treatment-effect object is
\[
G_{i2}(j,k)
=
Y_{i2}(k)-Y_{i2}(j).
\tag{K.9}
\]
It asks how worker \(i\)'s wage would change if firm \(j\) were replaced by
firm \(k\).  Subtracting (K.1) at the two firms gives
\[
\begin{aligned}
G_{i2}(j,k)
&=
\psi_k-\psi_j\\
&\quad+
\varpi_{ik}-\varpi_{ij}\\
&\quad+
\widetilde\varepsilon_{i2}(k)
-\widetilde\varepsilon_{i2}(j).
\end{aligned}
\tag{K.10}
\]
The common worker component \(\alpha_i\) and the worker-time covariate term
\(X_{i2}'\beta\) cancel exactly.

The wage object in (K.9) is closest to Sorkin--Warwar's formulation.
Sorkin--Warwar explicitly interpret a firm's effect as a worker-specific
treatment effect: exogenous mobility permits heterogeneous effects, whereas
additive separability makes the permanent treatment effect constant across
workers.  Volpe's displayed primitive is instead the wage level
\(Y_{it}\), with exogenous mobility and additive separability stated as
separate restrictions on its residual and functional form.  A wage
difference is an implication of Volpe's level equation, not the organizing
primitive.

Under pure additive separability, \(\varpi_{ij}=0\), so the permanent
conditional-mean part of (K.10) is the common firm contrast
\(\psi_k-\psi_j\).  Kline instead assumes
\[
\mathsf{NSG}_{K}:
\quad
G_{i2}(j,k)\perp\mathbf D_i
\quad\text{for every }j\neq k.
\tag{K.11}
\]
Condition (K.11) does not set
\(\varpi_{ik}-\varpi_{ij}\) to zero.  It permits heterogeneous gains but
rules out selection into the mobility path on those gains.  Consequently,
the wage object is closest to Sorkin's treatment-effect object, and Kline's
restriction on that object is conceptually closer to exogenous mobility
than to additive separability.

The restrictions are not literally equivalent.  Kline imposes full
independence on a difference, while Sorkin--Warwar impose conditional-mean
restrictions on wage-level components.  A worker component common to both
potential wages can be correlated with mobility and still cancel from
\(G_{i2}(j,k)\).  Conversely, conditional mean zero of the level components
does not generally imply full distributional independence of the gain.

\textit{Sources: Kline, pp.\ 32--33, Assumption 4 and Proposition 2;
Sorkin--Warwar, pp.\ 6--7.  Equations (K.9)--(K.10) and the comparison are
derived.}

\subsection{Why an additive average edge effect is not additive separability}

Under Kline's exclusion, parallel-trends, and stationarity assumptions, the
mean observed wage change on edge \(j\to k\) identifies the average
treatment effect for workers selected onto that edge:
\[
\Delta_{jk}
=
\mathbb E[
G_{i2}(j,k)
\mid D_{i1}=j,D_{i2}=k].
\tag{K.12}
\]
This causal edge result does not require additive separability.

If Kline's no-selection condition (K.11) is added, then
\[
\begin{aligned}
\Delta_{jk}
&=
\mathbb E[
G_{i2}(j,k)
\mid D_{i1}=j,D_{i2}=k]\\
&=
\mathbb E[G_{i2}(j,k)]\\
&=
\mathbb E[Y_{i2}(k)]
-\mathbb E[Y_{i2}(j)].
\end{aligned}
\tag{K.13}
\]
The first equality is the definition of the edge ATT.  The second uses
no selection on the gain.  The third uses linearity of expectation.
Define, relative to reference firm \(0\),
\[
\psi_j^{\mathrm{ATE}}
=
\mathbb E[Y_{i2}(j)]
-\mathbb E[Y_{i2}(0)].
\tag{K.14}
\]
Subtracting the two definitions in (K.14) gives
\[
\Delta_{jk}
=
\psi_k^{\mathrm{ATE}}
-\psi_j^{\mathrm{ATE}}.
\tag{K.15}
\]
Thus the average edge effects are transitive and admit a firm-difference
representation.

Equation (K.15) is an additive representation of population averages.  It
does not say that an individual's potential wage has the form
\(\alpha_i+\psi_j\), and it does not say that the individual gain equals
\(\psi_k^{\mathrm{ATE}}-\psi_j^{\mathrm{ATE}}\).  Therefore (K.15) is not
additive separability and is not equivalent to it.

\textit{Source: Kline, pp.\ 30--33, Propositions 1--2.  The line-by-line
expansion in (K.13)--(K.15) is derived.}

\subsection{A two-firm proof by counterexample}

Consider two firms, \(0\) and \(1\), set covariates and transitory errors to
zero, and let
\[
\varpi_{i0}=\xi_i,
\qquad
\varpi_{i1}=-\xi_i,
\qquad
\mathbb E[\xi_i]=0,
\tag{K.16}
\]
where \(\xi_i\) is nondegenerate and independent of the mobility path
\(\mathbf D_i\).  Potential wages are
\[
Y_{i2}(0)=\alpha_i+\psi_0+\xi_i,
\qquad
Y_{i2}(1)=\alpha_i+\psi_1-\xi_i.
\tag{K.17}
\]
Subtracting the two equations yields the individual gain
\[
G_{i2}(0,1)
=
\psi_1-\psi_0-2\xi_i.
\tag{K.18}
\]
Because \(\xi_i\) is nondegenerate, the gain varies across workers.
Moreover, \(\varpi_{i0}\) and \(\varpi_{i1}\) are not identically zero, so
pure additive separability \(\mathsf A_0\) fails.

Because \(\xi_i\perp\mathbf D_i\), equation (K.18) implies
\[
G_{i2}(0,1)\perp\mathbf D_i,
\tag{K.19}
\]
so Kline's no-selection condition holds.  Finally,
\[
\begin{aligned}
\Delta_{01}
&=
\mathbb E[G_{i2}(0,1)]\\
&=
\psi_1-\psi_0-2\mathbb E[\xi_i]\\
&=
\psi_1-\psi_0.
\end{aligned}
\tag{K.20}
\]
The average edge effect is an additive firm difference, while every
individual gain still contains the heterogeneous interaction
\(-2\xi_i\).  This proves by construction that an additive representation
of average edge effects does not imply individual-level additive
separability.

\textit{Equations (K.16)--(K.20) are a derived counterexample.}

\section*{References and page convention}

\begin{enumerate}
\item Borovi\v{c}kov\'a, Katar\'ina, and Robert Shimer.
\emph{High Wage Workers Work for High Wage Firms}, February 13, 2020.
Page references use the paper's printed page numbers.
\item Kline, Patrick M.  \emph{Firm Wage Effects}.  NBER Working Paper
33084, October 2024, revised February 2025.  Page references use the
paper's printed page numbers.
\item Sorkin, Isaac, and Lucas Warwar.  \emph{Paired Movers}, March 2026.
Page references use the paper's printed page numbers.
\item Sorkin, Isaac.  \emph{Labor Economics, Lecture 10}.  The cited slide
is numbered 46/79 in the lecture.
\item Volpe, Oscar.  \emph{Labor Omnibus Lectures, November 3 and 5}.
The cited slide is numbered 4/35 in the lecture.
\end{enumerate}

\end{document}
